% ------------------------------------------------------------------------
% file `3_uaa1_fonctions_2022-2023-connaitre-2-solution-body.tex'
%   in folder `xsim/'
%
%     solution of type `connaitre' with id `2'
%
% generated by the `solconnaitre' environment of the
%   `xsim' package v0.21 (2022/02/12)
% from source `3_uaa1_fonctions_2022-2023' on 2022/10/25 on line 53
% ------------------------------------------------------------------------
    1 point par justification correcte.
    \begin{enumerate}
        \item Faux.
              $\begin{aligned}[t]
                      14 & \stackrel{?}{=} 3\cdot 4+ 4 \\
                      14 & \stackrel{?}{=} 12 + 4      \\
                      14 & \ne 16
                  \end{aligned}$
        \item La racine d'une fonction est l'abscisse du point d'intersection de la fonction avec l'axe des abscisses.
        \item Faux. $f(x)= 8 = x^0 + 8$, c'est une fonction du degré 0.
        \item Vrai. Deux fonctions qui ont la même pente ont des graphes parallèles. Ici, dans les deux cas la pente vaut 3.
    \end{enumerate}
